\section{Stalks of abelian presheaves}
\label{section-stalks-abelian-presheaves}

\noindent
We first deal with the case of abelian groups as
a model for the general case.

\begin{lemma}
\label{lemma-stalk-abelian-presheaf}
Let $X$ be a topological space. Let $\mathcal{F}$ be a presheaf
of abelian groups on $X$. There exists a unique structure of an
abelian group on $\mathcal{F}_x$ such that for every
$U \subset X$ open, $x\in U$ the map $\mathcal{F}(U) \to \mathcal{F}_x$
is a group homomorphism. Moreover,
$$
\mathcal{F}_x
=
\colim_{x\in U} \mathcal{F}(U)
$$
holds in the category of abelian groups.
\end{lemma}

\begin{proof}
We define addition of a pair of elements
$(U, s)$ and $(V, t)$ as the pair $(U \cap V, s|_{U\cap V} +
t|_{U \cap V})$. The rest is easy to check.
\end{proof}

\noindent
What is crucial in the proof above is that the partially ordered set of open
neighbourhoods is a directed set (Categories, Definition
\ref{categories-definition-directed-set}). Namely, the coproduct
of two abelian groups $A, B$ is the direct sum $A \oplus B$, whereas
the coproduct in the category of sets is the disjoint
union $A \amalg B$, showing that colimits in the category
of abelian groups do not agree with colimits in the
category of sets in general.
